% Français 9115, batch 33 — Gallica views 641–660.
% Pass: Opus 5.5 (claude-opus-5-5), 2026-09-22 — first pass, unchecked against the views by a human.
%
% What the twenty views hold. Not number theory, and not the plates: the
% analysis of heat in a sphere and a cylinder, written out in spherical
% coordinates, with one leaf on the curvature of the sphere. Every leaf is
% mounted on a guard and written on one side only (except 316, written on
% both sides); nine views are blank backs or a blank leaf. Nine views are
% transcribed: 643, 644, 645, 646, 650, 652, 654, 656, 660. They are hers
% as far as the hand shows — the same regular hand as the plate memoirs,
% her spelling (« dela », « parconséquent », « complette », « tems »,
% « seroit », « voyoit », « spérique », « imployée »).
%
%   v. 643      (slip inked 315, the number on its blank side, v. 642)
%               calculation only: change of variables for a function v
%               from x, y, z to r, θ, φ about an axis inclined at A
%               (x = r sin A cosθ/sinθ cosφ, dz = sin A dr), the first
%               derivatives dr/dx … dθ/dz, dv/dx, dv/dy, dv/dz, and the
%               start of d²v/dx², whose last line runs under the guard.
%   v. 644      (leaf 316) scratch calculation: the second derivatives
%               again with θ′ (sin θ′ = cos θ) and r′, the terms in
%               d²v/dr², d²v/dθ², dv/dθ reduced, an « hypothese »
%               1/cos²θ circled; ends « Il faut encore ajouter a … » and a
%               crossed-out line; a pen trial and a large isolated 4.
%   v. 645      (back of 316) a short paragraph headed « cylindre »: if
%               the cylinder is heated arbitrarily, d²v/dz² cannot be
%               supposed zero; the equation (m/k)u + d²u/dr² + (1/r)du/dr
%               + (1/r²)d²u/dφ² + d²u/dz² = 0, and a particular solution
%               u = sin δφ r^β s sin(νz/ρ) with β = √(δ² + ν²/ρ²).
%               It cites « p. 3 » and « p. 4 ».
%   v. 646      (leaf 317) prose: the same equation reached by the motion
%               of heat through a molecule of the sphere bounded by two
%               cones, two planes and two spherical caps; sides r dθ,
%               r sinθ dφ, dr; faces; volume r² sinθ dθ dφ dr; the heat
%               accumulated in dt, K/r²(d(r² dv/dr)/dr + d(sinθ dv/dθ)/sinθ
%               dθ + 1/sin²θ d²v/dφ²) r² sinθ dθ dφ dr, divided by
%               c.D r² sinθ dθ dφ dr. « dv/dt = » is left without its right-
%               hand side. Pen trials (« Sphere ») in the margin.
%   v. 650      (leaf 319) « Somme des raisons inverses des deux rayons de
%               principale courbure »: 1/R + 1/R′ in dz/dx, dz/dy and the
%               second derivatives, then on the sphere (x = r sinθ cosφ,
%               y = r sinθ sinφ, r constant): dz/dx, dz/dy, K³, d²z/dxdy.
%               The one leaf of the batch not about heat.
%   v. 652      (leaf 320, marked « 1 » with a cross above the ink number)
%               the opening of a piece titled « Équation du mouvement varié
%               dela chaleur dans une sphère solide échauffée d'une manière
%               entièrement arbitraire »: dv/dt = K/c.D (d²v/dx² + d²v/dy² +
%               d²v/dz²) to be written in polar coordinates, « en adoptant
%               les dénominations qui appartiennent a la théorie dela
%               terre » (equator, pole, parallel); x, y, z and dx, dy, dz;
%               stops on « on tire ».
%   v. 654      (leaf 321, marked with a struck figure read 2) the
%               continuation: dx sinθ cosφ + … = dr, … = r dθ, … = r sinθ
%               dφ; the nine partial derivatives; dv/dx, dv/dy, dv/dz.
%   v. 656      (leaf 322, headed « Note ») from « l'équation (A) p. 3 dela
%               sphère », r = ∞, r sinθ = s: the sphere's equation becomes
%               the cylinder's. « cette maniere dela déduire est plus simple
%               que celle imployée par M^r Poisson car je ne vois pas la
%               necessité de mettre r − a a la place de r ». A paragraph
%               and equation crossed out with a large X; then why the
%               cylinder is contained in the sphere (the cone on the
%               spherical cap becomes a cylinder when the radius is
%               infinite).
%   v. 660      (leaf 324) the second derivatives d²v/dx², d²v/dy²,
%               d²v/dz² in r, θ, φ, their sum, and « L'équation générale
%               devient donc » dv/dt = K/c.D {2/r dv/dr + d²v/dr² +
%               d²v/r²dθ² + cosθ/r² sinθ dv/dθ + 1/r² sin²θ d²v/dφ²}
%               …. (A) — the equation (A) that v. 656 cites.
%
% The order of the pieces as bound is not the order of the argument:
% v. 652 → 654 → 660 is one continuous derivation ending on (A), while the
% « Note » that uses (A) (v. 656) and the blank leaf 323 (v. 658) are bound
% between 321 and 324; the prose of v. 646 and the « cylindre » paragraph
% of v. 645 refer to pages 3 and 4 not identified here. Recorded, not
% reordered.
%
% Views not transcribed. v. 641: blank back of the leaf inked 314 (v. 640,
% previous batch), its text showing through reversed. v. 642: the slip
% inked 315 seen from its blank side, the number written there, the
% calculation of v. 643 showing through reversed (mirrored and compared).
% v. 647, 649, 651, 653, 655, 657: blank backs of 317, 318, 319, 320, 321,
% 322, the recto showing through. v. 648: leaf inked 318, blank on both
% sides (checked at raised contrast). v. 658: leaf inked 323, blank.
% v. 659: its back, blank. No repeat, no béquet, no microfilm target.
%
% Seams. View 641 is a blank back; the last written view before the batch,
% v. 640 (ink 314, glanced at only), is a calculation in r, θ, φ, then
% « Recherches sur les cones » (batch 32, on disk only after this pass had
% read its views, reads « Recherches sur les \struck{surfa} cones »): the
% heat crossing the element of a cone, stopping three quarters down the
% leaf. Same apparatus as v. 643–646, so the batch continues a subject
% begun before it; but no sentence is cut at 640/641. View 660 ends
% complete on the equation (A), the foot of the leaf blank; v. 661 is its
% blank back. Nothing is cut at 660/661.
%
% The numbers on the leaves, and why no \folio{} is written. (1) The ink
% series at the top right, in the larger rounder hand: 315 (v. 642, on the
% slip's blank side), 316 (644), 317 (646), 318 (648, a blank leaf), 319
% (650), 320 (652), 321 (654), 322 (656, underlined), 323 (658, a blank
% leaf), 324 (660, a stroke under it). One number per leaf, blank leaves
% included, never on the backs, continuing 314 at v. 640 — penned, not
% pencilled, so, as everywhere in this volume, no \folio{}. (2) Marks above
% the ink number on two leaves, in a thicker pen: « 1 » followed by a small
% cross (v. 652, the leaf that opens the heat memoir) and a figure struck
% through with an oblique stroke, read 2 with doubt (v. 654, its sequel):
% perhaps her own page numbers for the piece; not settled. (3) « Note » at
% the head of v. 656. (4) Page references in the text: « p. 3 » (v. 645,
% 656), « p. 4 » (v. 645), none of which matches a number read on these
% leaves. (5) An isolated large « 4 » at the foot of v. 644, with a pen
% trial. Every number above was read on its view; none is inferred.
%
% Manuscript C. None of these leaves can be Laubenbacher and Pengelley's
% « Manuscript C » (ff. 348r–349r): the subject is heat and curvature, and
% the ink series here runs 315–324. At one number per leaf, 348 would fall
% some fifty views further on — an estimate, not a reading.
%
% Notation. Hers, kept. « sin », « cos », « tang » with or without a stop
% are set \text{sin}, \text{cos}, \text{tang}; θ, φ her angles, r the
% radius, A the fixed angle of v. 643, v the temperature. The capital in
% « K/c.D » is a dome on a bar, set D (as batch 28 set it at v. 549); it
% could be an Ω. The letter set ν at v. 645 is a looped v. Words she
% underlines in prose are set \emph{}. Braces opened with « { » and closed
% with « ) », or not closed, are kept as written. A long s is transcribed
% s. Passages cancelled by a large X are transcribed with a note, not
% inside \struck{}.
%
% What could not be read, and what is offered with doubt. \ill{} stands
% sixteen times: the blotted sign between two terms (v. 643); the
% exponent of cos in the first brace, the blotted leading signs of two
% calculations, two struck groups, the exponent of cos in the last line of
% the left block, the leading symbol before d²v/dθ² and a struck
% denominator in the last line (v. 644, eight in all); the sign between dφ
% and dr, a small blackened word, and a word among the marginal pen trials
% (v. 646); the value of cos θ under a blot (v. 656); a struck word in the
% first line of the third block (v. 654); the exponent of sin θ in the
% last brace of d²v/dy² (v. 660). \uncertain{} stands for the exponent 3
% (v. 643); « hypothese » (v. 644, first occurrence) and « Somme » (the
% vertical word, v. 644, in a note); « des » (v. 646); the factor 2 before
% cos²θ in K³ (v. 650); x, « ba » and c (v. 656); the coefficients 2 and 1
% and the letter A of the reference (v. 660). Nothing was guessed to fill
% a gap.
%
% Nothing here was checked against any published edition of Sophie
% Germain's papers or of any memoir on heat or on curvature.
\documentclass[11pt,a4paper]{article}
% The preamble shared by every transcription of the mathematicians' manuscripts in Gallica.
%
% It rests on one idea: **uncertainty compounds, it does not resolve.** A
% transcription that smooths over a doubtful word has destroyed the only thing
% it was made for — a reader must be able to tell, at every point, what was on
% the page and what was supplied. The apparatus macros below are therefore the
% critical apparatus, not decoration.
%
% The same commands are recognised by `scripts/render.mjs`, which produces the
% site's reading view, and by `scripts/tei.mjs`. A macro added here must be
% added there too, or rendering fails loudly — which is the wanted behaviour.
%
% Comments here are English, like the rest of the repository. The documents
% this preamble sets are French, and that is deliberate: see the skills.

\ProvidesPackage{germain}[2026/09/15 Transcriptions of mathematicians' manuscripts in Gallica]

\usepackage{amsmath,amssymb,amsthm}
\usepackage{mathrsfs}
% Kept from the parent preamble so the renderer's diagram support stays
% available; these manuscripts draw no commutative diagrams, and nothing here uses it.
\usepackage{tikz-cd}
\usepackage[english,french]{babel}
\usepackage{fontspec}

% --- Greek ------------------------------------------------------------------
% The text face stays fontspec's default, Latin Modern, which has no Greek:
% XeTeX drops every glyph it lacks with a line in the log and nothing on the
% page, and the Greek words the manuscripts quote vanished from the PDFs. So
% the Greek and Greek Extended blocks get a XeTeX character class of their own,
% and entering or leaving a run of it switches the family to CMU Serif — the
% same Computer Modern design, with polytonic Greek — and back. Only the
% family changes: a Greek word in \textit or \textbf keeps its shape and series.
% The transitions to and from every other class carry the tokens that class
% has towards an ordinary letter, so French spacing before « ; » or inside
% « » still holds after a Greek word. No grouping, so no brace or \endgroup
% can come out unbalanced; maths is left alone, since XeTeX inserts nothing
% there. Kept identical in `exercices.sty`.
\makeatletter
\newfontfamily\germainGreekfont{cmun}[
  Extension=.otf, NFSSFamily=germaingreek,
  UprightFont=*rm, ItalicFont=*ti, SlantedFont=*sl,
  BoldFont=*bx, BoldItalicFont=*bi, BoldSlantedFont=*bl]
\newXeTeXintercharclass\germain@greek
\def\germain@greekrange#1#2{%
  \@tempcnta=#1\relax
  \loop
    \XeTeXcharclass\@tempcnta=\germain@greek
    \ifnum\@tempcnta<#2\advance\@tempcnta\@ne\repeat}
\germain@greekrange{"0370}{"03FF}
\germain@greekrange{"1F00}{"1FFF}
\def\germain@greekprev{\rmdefault}
\def\germain@greekin{%
  \edef\germain@greekprev{\f@family}\fontfamily{germaingreek}\selectfont}
\def\germain@greekout{\fontfamily{\germain@greekprev}\selectfont}
\def\germain@greekjoin{%
  \XeTeXinterchartokenstate=\@ne
  \@tempcnta=\z@
  \loop
    \ifnum\@tempcnta=\germain@greek\else
      \XeTeXinterchartoks\germain@greek\@tempcnta=\expandafter{%
        \expandafter\germain@greekout\the\XeTeXinterchartoks\z@\@tempcnta}%
      \XeTeXinterchartoks\@tempcnta\germain@greek=\expandafter{%
        \the\XeTeXinterchartoks\@tempcnta\z@\germain@greekin}%
    \fi
    \ifnum\@tempcnta<4095\advance\@tempcnta\@ne\repeat}
% babel-french sets its own classes, and saves and restores the interchar
% state, each time a language is selected: join again after it does.
\addto\extrasfrench{\germain@greekjoin}
\addto\extrasenglish{\germain@greekjoin}
\AtBeginDocument{\germain@greekjoin}
\makeatother

\usepackage{geometry}
\geometry{a4paper,margin=25mm,marginparwidth=28mm}
\usepackage{marginnote}
\usepackage{xcolor}
\usepackage[normalem]{ulem}
\setlength{\emergencystretch}{3em}
\usepackage{hyperref}
\hypersetup{colorlinks=true,linkcolor=black,urlcolor=black,pdfborder={0 0 0}}

% --- Batch metadata ---------------------------------------------------------
% Taken as they stand by the rendering script, which shows them at the head of
% the reading view. They fix what the file claims to cover. The macro names are
% the parent project's, so that the scripts stay shared; \folder here means the
% volume's slug — `fr-9115` — and \pages counts Gallica views.

\newcommand{\folder}[1]{\gdef\grothFolder{#1}}
\newcommand{\batch}[1]{\gdef\grothBatch{#1}}
\newcommand{\pages}[2]{\gdef\grothFirst{#1}\gdef\grothLast{#2}}
\newcommand{\foldertitle}[1]{\gdef\grothFoldertitle{#1}}

% The holder's own shelfmark, printed as they write it: « Français 9115 ».
\newcommand{\shelfmark}[1]{\gdef\grothShelfmark{#1}}

% The dating, copied verbatim from the holder's catalogue — never inferred
% here. « XIXe siècle » is the BnF's; a pass that dates a leaf from its content
% says so in a \note{}, as its own inference.
\newcommand{\dating}[1]{\gdef\grothDating{#1}}

% The status notice, printed in the title block and repeated as a diagonal
% watermark on every page of the PDF. The manuscripts are in the public
% domain; what this line declares is not a right but a quality — an
% unverified first pass by a machine. The skills write the line; a file
% without it gets no watermark, so leaving it out is visible in the output.
\newcommand{\watermark}[1]{\gdef\grothWatermark{#1}}

\gdef\grothFolder{?}\gdef\grothBatch{}\gdef\grothFirst{?}\gdef\grothLast{?}%
\gdef\grothFoldertitle{}\gdef\grothShelfmark{}\gdef\grothDating{}\gdef\grothWatermark{}

% --- The résumé -------------------------------------------------------------
\newenvironment{resume}
  {\small\setlength{\parskip}{0.45em}}
  {\par\medskip\hrule\medskip}

% --- Keywords ---------------------------------------------------------------
% In English, closing the modernised reading's résumé: the modern vocabulary
% under which to search for what the leaves construct. The manifest extracts
% them to tag the volume — this line is the single source of the tags.
\newcommand{\keywords}[1]{\par\smallskip\noindent
  {\footnotesize\textsc{Keywords} --- \textit{#1}}\par}

% --- The critical apparatus -------------------------------------------------

% The Gallica view number, in the margin. This is the anchor that lets a
% disagreement over a formula be settled by looking at *one* image rather than
% a volume of 750. The site uses it to turn the facsimile as the transcription
% is scrolled. A view is one scanned image: usually one side of a leaf.
\newcommand{\page}[1]{%
  \marginnote{\footnotesize\color{gray!70}\textsf{v.\,#1}}%
  \label{page:#1}\ignorespaces}

% The BnF's pencilled foliation, where the leaf carries it and a pass can read
% it: « 198r ». This is what the literature cites — « ff. 198r–208v » — and
% the one line that turns a citation into a view. Recorded, never inferred:
% a folio a pass cannot read is not written.
\newcommand{\folio}[1]{%
  \marginnote{\footnotesize\color{gray!50}\textsf{f.\,#1}}\ignorespaces}

% The modernised reading's coarser counterpart to \page{}: which views a
% section is drawn from.
\newcommand{\pagerange}[2]{%
  \marginnote{\footnotesize\color{gray!70}\textsf{v.\,#1--#2}}%
  \ignorespaces}

% Illegible: never guessed.
\newcommand{\ill}{\textcolor{red!60!black}{[\,\ldots\,]}}

% A doubtful reading: the word is given, and flagged as uncertain.
\newcommand{\uncertain}[1]{\uline{#1}}

% An editorial addition — a missing letter, an implied word.
\newcommand{\add}[1]{\textcolor{blue!50!black}{[#1]}}

% Struck out by the writer. What was crossed out is part of the document.
\newcommand{\struck}[1]{\sout{#1}}

% The transcriber's note, on the state of the page or the mathematics.
\newcommand{\note}[1]{\par\smallskip{\small\color{gray!60!black}\textit{[#1]}}\par\smallskip}

% A marginal note of hers, distinct from the body of the text.
\newcommand{\marginal}[1]{\par\smallskip\noindent\hspace{1em}%
  {\small\color{gray!60!black}#1}\par\smallskip}

% --- Title ------------------------------------------------------------------
% The title block is French, like the documents themselves.

\AddToHook{shipout/background}{%
  \ifx\grothWatermark\empty\else
    \begin{tikzpicture}[remember picture, overlay]
      \node[rotate=52, scale=3.4, align=center, text=gray!18,
            font=\sffamily\bfseries]
        at (current page.center) {\grothWatermark};
    \end{tikzpicture}%
  \fi}

\AtBeginDocument{%
  \begin{center}
    {\large\bfseries Manuscrits de mathématiciens (Gallica) ---
      \ifx\grothShelfmark\empty\grothFolder\else\grothShelfmark\fi}\\[2pt]
    {\itshape\grothFoldertitle}\\[4pt]
    {\small vues \grothFirst--\grothLast
      \ifx\grothBatch\empty\else\ (lot \grothBatch)\fi}\\[2pt]
    \ifx\grothDating\empty\else
      {\small\color{gray!60!black}Datation du catalogue : \grothDating}\\[2pt]
    \fi
    \ifx\grothWatermark\empty\else
      {\small\bfseries\color{red!45!gray}\grothWatermark}\\[2pt]
    \fi
    {\footnotesize\color{gray!60!black}Transcription établie d'après les images IIIF de Gallica.
     Source gallica.bnf.fr / Bibliothèque nationale de France.\\
     Les lectures incertaines sont soulignées, les illisibles notées [\,\ldots\,],
     les ajouts entre crochets ; \textsf{v.} numérote les vues de Gallica,
     \textsf{f.} la foliotation au crayon de la BnF.}
  \end{center}
  \vspace{1em}}

\endinput


\folder{fr-9115}
\shelfmark{Français 9115}
\batch{33}
\pages{641}{660}
\dating{1801-1900}
\watermark{Lecture automatique\\non vérifiée}
\foldertitle{Papiers de Sophie GERMAIN. — Recueil de dissertations et problèmes mathématiques et physiques.}

\begin{document}

\note{Numérotation des feuillets. En haut à droite, à l'encre, d'une écriture plus ronde, la série qui court dans tout le volume : 315 (vue 642, porté sur le côté blanc du petit feuillet dont la vue 643 montre l'autre côté), 316 (vue 644), 317 (vue 646), 318 (vue 648, feuillet blanc), 319 (vue 650), 320 (vue 652), 321 (vue 654), 322 (vue 656), 323 (vue 658, feuillet blanc), 324 (vue 660) ; jamais au dos. Cette série numérote des feuillets, feuillets blancs compris ; elle occupe la place de la foliotation de la Bibliothèque, mais elle est tracée à la plume et non au crayon, et aucun « f. » n'est donc porté. Au-dessus du nombre, d'une plume plus épaisse : « 1 » suivi d'une petite croix (vue 652) et un chiffre barré, lu 2 avec doute (vue 654) ; « Note » en tête de la vue 656. Tous ces nombres ont été lus sur la vue ; aucun n'est déduit. Ne sont pas transcrites : la vue 641, dos blanc du feuillet 314 ; la vue 642, côté blanc du petit feuillet 315 ; les vues 647, 649, 651, 653, 655 et 657, dos blancs ; les vues 648 et 659, feuillet 318 et dos du feuillet 323, blancs ; la vue 658, feuillet 323, blanc.}

\page{643}

\note{Petit feuillet monté sur une feuille de garde, écrit de ce côté seulement ; le nombre « 315 », à l'encre, est porté sur l'autre côté, blanc (vue 642), où ce calcul se lit à l'envers par transparence. En haut de la vue, au-dessus du feuillet, le dos de la feuille 314 (vue 640) laisse voir son texte à l'envers. Le bord droit du feuillet passe sous la garde suivante et cache la fin de la dernière ligne. Calcul sans texte suivi.}

\[ \frac{dv}{dx} = \frac{dv}{d\theta}.\frac{d\theta}{dx} + \frac{dv}{d\varphi}.\frac{d\varphi}{dx} \]
\[ = \frac{\text{cos}\,\varphi\,\text{cos}^{2}(A-\theta)}{r\,\text{cos}A}.\frac{dv}{d\theta}\ \ill{}\ \frac{\text{sin.}\,\varphi}{r\,\text{cos.}A\,\text{tang}(A-\theta)}\,\frac{dv}{d\varphi} \]
\note{Le signe entre les deux termes est couvert d'une tache d'encre.}

\[ \frac{dv}{dy} = \frac{dv}{d\theta}.\frac{d\theta}{dy} + \frac{dy}{d\varphi}.\frac{d\varphi}{dy} \]
\[ = \frac{\text{sin}\,\varphi.\text{cos}^{2}(A-\theta)}{r\,\text{cos.}A}\,\frac{dv}{d\theta} + \frac{\text{cos}\,\varphi}{r\,\text{cos}A\,\text{tang}(A-\theta)}\,\frac{dv}{d\varphi} \]
\note{« $\frac{dy}{d\varphi}$ » est écrit ainsi, là où l'on attend $\frac{dv}{d\varphi}$. Le $\theta$ de « tang$(A-\theta)$ » est récrit et empâté.}

\[ dz = \text{sin.}A\,dr \qquad x = r\,\text{sin}A\,\frac{\text{cos}\,\theta}{\text{sin}\,\theta}\,\text{cos}\,\varphi \qquad y = r\,\text{sin.}A\,\frac{\text{cos.}\theta}{\text{sin}\,\theta}\,\text{sin.}\,\varphi \]
\note{« $dz = \text{sin.}A\,dr$ » est écrit dans la marge de gauche, en face des expressions de $x$ et de $y$.}
\[ dx = \text{sin.}A\,\frac{\text{cos}\,\theta}{\text{sin}\,\theta}\,\text{cos}\,\varphi\,dr + r\,\text{sin}A\,\frac{\text{cos}\,\theta}{\text{sin}\,\theta}\,\text{sin.}\,\varphi\,d\varphi - r\,\text{sin}A\,\text{cos}\,\varphi\,\frac{d\theta}{\text{sin}^{2}\theta} \]
\[ dy = \text{sin.}A\,\frac{\text{cos}\,\theta}{\text{sin}\,\theta}\,\text{sin.}\,\varphi\,dr - r\,\text{sin.}A\,\frac{\text{cos}\,\theta}{\text{sin.}\theta}\,\text{cos}\,\varphi\,d\varphi - r\,\text{sin.}A\,\text{sin.}\varphi\,\frac{d\theta}{\text{sin}^{2}\theta} \]
\note{Le signe $+$ devant le deuxième terme de $dx$ est lu tel qu'il est écrit ; celui de $dy$ est $-$.}

\[ dx\,\text{cos.}\varphi + dy\,\text{sin}\,\varphi - dz\,\frac{\text{cos}\,\theta}{\text{sin.}\theta} = -r\,\text{sin}A\,\frac{d\theta}{\text{sin}^{2}\theta} \]
\[ dx\,\text{sin.}\varphi - dy\,\text{cos.}\varphi = r\,\text{sin.}A\,\frac{\text{cos}\,\theta}{\text{sin}\,\theta}\,d\varphi \]
\note{Le numérateur « cos$\,\theta$ » du terme en $dz$ est récrit sur un premier mot.}

\[ \frac{dr}{dz} = \frac{1}{\text{sin.}A} \qquad \frac{d\theta}{dx} = -\frac{\text{cos}\,\varphi\,\text{sin}^{2}\theta}{r\,\text{sin.}A} \qquad \frac{d\varphi}{dx} = \frac{\text{sin}\,\varphi\,\text{sin}\,\theta}{r\,\text{sin}A\,\text{cos}\,\theta} \]
\[ \frac{dr}{dx} = 0 \qquad \frac{d\theta}{dy} = -\frac{\text{sin.}\varphi\,\text{sin}^{2}\theta}{r\,\text{sin.}A} \qquad \frac{d\varphi}{dy} = -\frac{\text{cos}\,\varphi\,\text{sin}\,\theta}{r\,\text{sin}A\,\text{cos}\,\theta} \]
\[ \frac{dr}{dy} = 0 \qquad \frac{d\theta}{dz} = \frac{\text{sin}\,\theta\,\text{cos}\,\theta}{r\,\text{sin.}A} \]
\note{Le « $r$ » des numérateurs $\frac{dr}{dx}$, $\frac{dr}{dy}$ est récrit sur une autre lettre ; le « cos » de $\frac{d\theta}{dx}$ est repassé.}

\[ \frac{dv}{dx} = \frac{\text{sin}\,\varphi.\text{sin}\,\theta}{r\,\text{sin}A\,\text{cos}\,\theta}\,\frac{dv}{d\varphi} - \frac{\text{cos}\,\varphi\,\text{sin}^{2}\theta}{r\,\text{sin.}A}\,\frac{dv}{d\theta} \]
\[ \frac{dv}{dy} = -\frac{\text{cos}\,\varphi\,\text{sin}\,\theta}{r\,\text{sin.}A\,\text{cos.}\theta}\,\frac{dv}{d\varphi} - \frac{\text{sin.}\varphi\,\text{sin}^{2}\theta}{r\,\text{sin.}A}\,\frac{dv}{d\theta} \]
\[ \frac{dv}{dz} = \frac{1}{\text{sin}A}.\frac{dv}{dr} + \frac{\text{sin}\,\theta\,\text{cos}\,\theta}{r\,\text{sin}A}.\frac{dv}{d\theta} \]
\note{Le $v$ de $\frac{dv}{dr}$ est récrit sur une autre lettre.}

\[ \frac{d^{2}v}{dx^{2}} = \frac{dv}{d\varphi}\left\{\frac{\text{cos}\,\varphi.\text{sin}^{2}\theta\,\text{sin.}\varphi}{r^{2}\,\text{sin.}^{2}A\,\text{cos.}^{\uncertain{3}}\theta} - \frac{\text{cos}\,\varphi\,\text{sin}\,\varphi\,\text{sin}^{2}\theta}{r^{2}\,\text{sin}^{2}A\,\text{cos}^{2}\theta}\right) \]
\[ - \frac{dv}{d\theta}\left\{\frac{\text{sin}^{2}\varphi\,\text{sin}^{3}\theta}{r^{2}\,\text{sin}^{2}A\,\text{cos.}\theta} - \frac{2\,\text{cos}^{2}\varphi\,\text{sin}^{3}\theta\,\text{cos.}\theta}{r^{2}\,\text{sin}^{2}A}\right\} \]
\[ + \frac{d^{2}v}{d\varphi^{2}}\,\frac{\text{sin}^{2}\varphi\,\text{sin}^{4}\theta}{r^{2}\,\text{sin}^{2}A\,\text{cos}^{2}\theta} - \frac{d^{2}v}{d\varphi\,d\theta}.\frac{\text{sin}\,\varphi\,\text{cos}\,\varphi\,\text{sin}^{3}\theta}{r^{2}\,\text{sin}^{2}A\,\text{cos.}\theta} \]
\[ - \frac{d^{2}v}{d\varphi\,d\theta}.\frac{\text{sin}\,\varphi\,\text{cos.}\varphi\,\text{sin}^{3}\theta}{r^{2}\,\text{sin}^{2}A\,\text{cos.}\theta} + \frac{d^{2}v}{d\theta^{2}}\,\frac{\text{cos}^{2}\varphi\,\text{sin}\ldots}{r^{2}\,\text{sin}^{2}A} \]
\note{L'accolade ouverte après $\frac{dv}{d\varphi}$ se ferme par une parenthèse. L'exposant de « cos.$\theta$ » au premier dénominateur est lu 3 avec doute (2 est possible). Le « 3 » de $\text{sin}^{3}\theta$ dans le premier terme de la ligne suivante est récrit. La ligne se perd sous la garde : « cos$^{2}\varphi$ sin » puis l'exposant et la suite sont cachés ; les points sont de l'éditeur. Le bas du feuillet est blanc.}

\page{644}

\note{En haut à droite, à l'encre, « 316 ». Feuille de calcul, sans texte suivi ; le dos (vue 645) porte en tête un court texte sur le cylindre. Le changement de variable de la vue 643 est poursuivi ici avec un angle $\theta'$ et un rayon $r'$ ; les accents sont portés comme on les lit, et l'on ne peut pas toujours dire s'il en manque un.}

\[ \frac{d^{2}v}{dr\,d\theta} = \frac{d^{2}v}{dr'd\theta}.\frac{\text{sin.}A}{\text{sin.}\theta} + \frac{d^{2}v}{d\theta^{2}}.\frac{\text{sin.}A}{r'\,\text{cos.}\theta} \]
\[ - \frac{dv}{dr'}\,\frac{\text{sin}A\,\text{cos}\,\theta}{\text{sin}^{2}\theta} + \frac{dv}{d\theta}.\frac{\text{sin.}A\,\text{cos}\,\theta}{r'\,\text{sin.}^{2}\theta} \]
\note{Le premier membre $\frac{d^{2}v}{dr\,d\theta}$ est d'une encre plus pâle que le reste. Le signe devant $\frac{dv}{dr'}$ est une tache, avec un « $-$ » écrit au-dessous.}

\[ \text{sin}\,\theta' = \text{cos}\,\theta \qquad d\theta'\,\text{cos}\,\theta' = -\text{sin}\,\theta\,d\theta \]
\note{Ces deux égalités sont écrites à droite, en face des lignes qui suivent.}

\[ \frac{d^{2}v}{dr\,d\theta} = \frac{d^{2}v}{dr'd\theta'}.\frac{d\theta'}{d\theta}\,\frac{\text{sin}A}{\text{sin}\,\theta'} + \frac{d^{2}v}{d\theta'^{2}}.\frac{d\theta'}{d\theta}.\frac{\text{sin}A}{r'\,\text{cos}\,\theta'} \qquad -\frac{d\theta'}{d\theta} = \frac{\text{sin}\,\theta}{\text{cos.}\theta'} \]
\[ - \frac{dv}{dr'}.\frac{\text{sin}A\,\text{cos}\,\theta'}{\text{sin}^{2}\theta'}\,\frac{d\theta'}{d\theta} + \frac{dv}{d\theta'}\,\frac{\text{sin}A\,\text{cos}\,\theta'}{r'\,\text{sin}^{2}\theta'}\,\frac{d\theta'}{d\theta} \]
\note{Le signe $-$ du début de la dernière ligne est écrit sur une tache.}

\[ \frac{1}{\text{sin}^{2}A}\,\frac{d^{2}v}{dr^{2}} + \frac{2\,\text{sin.}\theta\,\text{cos.}\theta}{r\,\text{sin}^{2}A}.\frac{d^{2}v}{dr\,d\theta} + \left\{\frac{\text{sin.}\theta\,\text{cos}\ill{}\theta}{r^{2}\,\text{sin}^{2}A} - \frac{\text{sin}\,\theta}{r\,\text{sin}^{2}A\,\text{cos.}\theta}\right\}\frac{dv}{d\theta} \]
\note{L'exposant de « cos » au premier terme de l'accolade est couvert d'une tache. Un petit « $-$ » isolé est écrit sous cette ligne, sous le premier dénominateur de l'accolade.}

\[ = \frac{1}{\text{sin}^{2}A}\left\{\frac{\text{sin}^{2}A}{\text{sin}^{2}\theta}\,\frac{d^{2}v}{dr'^{2}} + \frac{\text{sin}^{2}A}{r'\,\text{sin.}\theta\,\text{cos}\,\theta}\,2.\frac{d^{2}v}{dr'd\theta} + \frac{\text{sin}^{2}A}{r'^{2}\,\text{cos}^{2}\theta}\,\frac{d^{2}v}{d\theta^{2}} - \frac{\text{sin}^{2}A}{r'\,\text{sin}^{2}\theta}.\frac{dv}{dr'}\right. \]
\[ \left. - \left(\frac{\text{sin}^{2}A}{r'\,\text{cos}\,\theta\,\text{sin}\,\theta} - \frac{\text{sin}^{2}A}{r'^{2}\,\text{cos}^{3}\theta}\right)\frac{dv}{d\theta}\right\} \]
\note{Le dénominateur de $\frac{d^{2}v}{dr'^{2}}$ finit sur une tache ; celui de $\frac{d^{2}v}{d\theta^{2}}$ est récrit ; l'accent de $dv/dr'$ est surchargé. La parenthèse du dernier terme est écrite sur deux lignes, le second terme au-dessous du premier.}

\[ - \frac{2\,\text{sin}\,\theta\,\text{cos}\,\theta}{r'\,\text{sin}A\,\text{sin}\,\theta}\left\{\frac{\text{sin.}A}{\text{sin.}\theta'}.\frac{d^{2}v}{dr'd\theta'} + \frac{\text{sin.}A}{r'\,\text{cos.}\theta'}.\frac{d^{2}v}{d\theta'^{2}}\right\} \]
\[ \left(\frac{1}{\text{cos}^{2}\theta}\ \text{\uncertain{hypothese}}\right) \qquad \left(\text{acause de}\ \frac{d\theta'}{d\theta} = -1\right) \]
\note{Au dénominateur, après « sin$A$ », un mot récrit, lu « sin$\,\theta$ » sur un premier « cos ». Le « $+$ » entre les deux termes est écrit sur une tache. Les deux parenthèses sont écrites l'une au-dessus de l'autre à droite de l'accolade ; le mot de la première, d'une écriture rapide, est lu avec doute d'après le « hypothese » bien formé, et entouré, de la ligne plus bas ; « acause » en un mot.}

\[ = \frac{1}{\text{sin}^{2}\theta}.\frac{d^{2}v}{dr^{2}} + \frac{1}{r'^{2}\,\text{cos}^{2}\theta}.\frac{d^{2}v}{d\theta^{2}} + \frac{1}{r'\,\text{sin}^{2}\theta}.\frac{dv}{dr'} + \left(\frac{1}{r'^{2}\,\text{cos}\,\theta\,\text{sin}\,\theta} - \frac{\text{sin.}\theta}{r'^{2}\,\text{cos}^{3}\theta}\right)\frac{dv}{d\theta} \]
\note{À droite de cette ligne, un mot très petit, sans doute « hypothese » encore. Des traits courbes relient les termes de cette ligne au calcul de la moitié inférieure droite.}

\[ \left[\ill{}\left\{\frac{\text{cos}^{2}\theta - \text{sin}^{2}\theta}{r'^{2}\,\text{cos}^{3}\theta\,\text{sin}\,\theta}\right\}\frac{dv}{d\theta}\,\text{cos}^{2}\theta\,\text{sin}^{2}\theta\right] \]
\[ \frac{dv}{d\theta}\left(\frac{1}{\text{cos}^{2}\theta}\ \text{hypothese}\right).\frac{\text{sin}\,\theta\,\text{cos}\,\theta}{r'^{2}} \]
\[ = \frac{dv}{d\theta}.\frac{\text{sin}^{2}\theta\,\text{cos}^{2}\theta}{r'^{2}\,\text{cos}^{3}\theta\,\text{sin.}\theta}.\frac{1}{\text{cos}^{2}\theta} \]
\[ = \frac{1}{r'^{2}\,\text{cos}^{\ill{}}\theta\,\text{sin}\,\theta}\,\frac{dv}{d\theta} \]
\note{Ces quatre lignes, à gauche, sont enfermées dans une grande accolade tracée à la main. Le signe du début est une tache. « $\frac{1}{\text{cos}^{2}\theta}$ hypothese » est entouré d'un trait. L'exposant de « cos » à la dernière ligne est écrit sur une tache.}

\[ \ill{}\left\{\frac{\text{cos}\,\theta\,\text{sin}^{2}\theta}{r'^{2}\,\text{sin}^{2}\theta} - \frac{\text{cos}\,\theta}{r'\,\text{sin}^{3}\theta}\right\}\frac{dv}{d\theta}.\frac{\text{sin}^{2}\theta}{\text{cos}^{2}\theta} \]
\[ = \frac{dv}{d\theta}\left\{\frac{\struck{\ill{}}\ \text{sin}\,\theta\,\text{cos}\,\theta\,(\text{sin}^{2}\theta + \text{cos}^{2}\theta)}{r'^{2}\,\text{cos}^{4}\theta}\ \struck{\ill{}}\right) \]
\[ - \frac{\text{sin}\,\theta\,\text{cos}\,\theta\,\text{cos}^{2}\theta}{r'^{2}\,\text{cos}^{4}\theta} \]
\[ + \frac{\text{cos}^{2}\theta}{r'^{2}\,\text{sin}^{3}\theta}.\frac{\text{sin}^{2}\theta}{\text{cos}^{2}\theta} \]
\note{À droite, la même réduction sur l'autre terme. Le signe du début est une tache. « sin$\,\theta$ cos$\,\theta$ » est écrit au-dessus d'un premier facteur noirci ; à droite de la deuxième ligne, un groupe entièrement noirci, où se devinent « sin », « cos », puis une parenthèse. Entre les deux colonnes, écrit de bas en haut, un mot lu \uncertain{Somme}, et au-dessous un petit signe.}

Il faut encore ajouter a $\frac{\text{sin}^{2}\theta\ \struck{\text{cos}}}{r'^{2}\,\text{cos}^{2}\theta\,\text{sin}^{2}\theta}\,\frac{d^{2}v}{d\theta^{2}}$
\[ \ill{}\ \frac{d^{2}v}{d\theta^{2}} = \frac{dv}{d\theta}\left\{\frac{\text{sin}^{2}\theta}{\text{cos}\ \struck{\ill{}}} - \frac{\text{cos}\,\theta}{r'^{2}\,\text{sin}\,\theta} + \frac{\text{sin}\,\theta}{r'^{2}\,\text{cos}^{3}\theta}\right. \]
\[ = \frac{dv}{d\theta}\left(\right. \]
\note{Sous « a », au-dessous de la fraction, « cos$^{2}\theta$ » souligné, d'un trait séparé. La ligne suivante commence, à gauche, par une fraction noircie ; son premier terme entre accolades est raturé de boucles serrées ; un long trait courbe barre toute la ligne, et « $= \frac{dv}{d\theta}($ » au-dessous, laissé sans suite. Plus bas, un essai de plume (« emme ») et un grand chiffre isolé, lu 4 ; le bas de la feuille est blanc, avec le texte de l'autre côté visible à l'envers.}

\page{645}

\note{Le dos de la feuille 316 (vue 644), sans numéro : un court texte en haut de la page ; le reste est blanc, le calcul de l'autre côté se lisant à l'envers par transparence. Le mot « cylindre » est en titre, souligné.}

\emph{cylindre}

Si on veut que le cylindre ait été échauffé arbitrairement
on ne pourra plus supposer $\frac{d^{2}v}{dz^{2}} = 0$ l'équation
transformée contiendra ce terme sans aucune altération
ensorte qu'elle deviendra p. 3
\note{« entierement » est écrit au-dessus de « arbitrairement », qui n'est pas barré. « p. 3 » et, plus bas, « p. 4 » renvoient à des pages qui ne sont pas identifiées dans ce lot.}
\[ \frac{m}{k}u + \frac{d^{2}u}{dr^{2}} + \frac{1}{r}.\frac{du}{dr} + \frac{1}{r^{2}}.\frac{d^{2}u}{d\varphi^{2}} + \frac{d^{2}u}{dz^{2}} = 0 \]
si aulieu deprendre la valeur de \emph{u} p. 4 on fait
$u = \text{sin}\,\delta\varphi\,r^{\beta}s\,\text{sin.}\frac{\nu z}{\rho}$ et qu'on prenne
$\beta = \sqrt{\left(\delta^{2} + \frac{\nu^{2}}{\rho^{2}}\right)}$ on satisfera a l'équation
\note{La lettre lue $\nu$ est une lettre bouclée comme un « v » ; celle lue $\rho$, au dénominateur, est suivie d'une tache, et l'exposant de $\nu$ sous le radical est récrit sur une tache. Le « $u$ » de « la valeur de $u$ » est souligné. La phrase s'arrête là ; le reste de la page est blanc.}

\page{646}

\note{En haut à droite, à l'encre, « 317 ». Feuillet monté sur une feuille de garde, écrit d'un seul côté, d'une écriture régulière mais corrigée entre les lignes ; le dos (vue 647) est blanc. Pas de titre ni de numéro de page.}

On parviendra a la même équation en considerant le mouvement
dela chaleur dans une molécule dela sphère cette molécule
sera comprise entre les deux surfaces coniques dont les côtés
font avec l'axe des $z$ qui est en même tems celui \uncertain{des}
\struck{deux} solides coniques. les angles $\theta$ et $\theta + d\theta$ entre les deux
plans qui font avec celui de $x$ et $z$ les angles $\varphi$ et $\varphi + d\varphi$
et enfin entre les deux calottes sphériques qui appartiennent
aux sphères dont les rayons sont $r$ et $r + dr$
\note{« molécule », dans la deuxième ligne, est récrit sur un premier mot noirci dont on voit un « p » ou un « s » initial ; « une » est répété dans l'interligne, au-dessous. À la fin de la quatrième ligne, « des » suivi d'un trait, la lecture du dernier signe incertaine. Au-dessus de « les angles $\theta$ et $\theta + d\theta$ », dans l'interligne, sans signe d'insertion : « auxquels ces surfaces appartiennent ».}

parconséquent \struck{$r\,d\theta$, $r$} $r\,d\theta$, $r\,\text{sin}\,\theta\,d\varphi$ \ill{} $dr$ \struck{seront} représenteront trois
des côtés, $r^{2}\,\text{sin}\,\theta\,d\theta\,d\varphi$, \struck{$r\,\text{sin}\,\varphi$} $r\,d\theta\,dr$, \struck{\ill{}} $r\,\text{sin}\,\theta\,d\varphi\,dr$,
trois des faces et $r^{2}\,\text{sin}\,\theta\,d\theta\,d\varphi\,dr$ le volume dela même
molécule.
\note{« représenteront » est écrit au-dessus de « seront », barré. Entre $d\varphi$ et $dr$ de la première ligne, une tache couvre un signe (une virgule ou « et ») ; devant $r\,\text{sin}\,\theta\,d\varphi\,dr$, un petit mot noirci.}

En retranchant les quantités de chaleurs qui entrent dans
la molécule, par chacune de ses trois premières faces, de celle qui s'échappe par les trois
\emph{dernieres} faces respectivement opposées \struck{aux premieres}, on voit qu'il s'accumule
durant l'instant $dt$ dans l'intérieur de cette molécule une
quantité de chaleur égale a la quantité
\[ \frac{K}{r^{2}}.\left(\frac{d\left(r^{2}.\frac{dv}{dr}\right)}{dr} + \frac{d\left(\text{sin}\,\theta.\frac{dv}{d\theta}\right)}{\text{sin}\,\theta\,d\theta} + \frac{1}{\text{sin}^{2}\theta}.\frac{d^{2}v}{d\varphi^{2}}\right)r^{2}\,\text{sin}\,\theta\,d\theta\,d\varphi\,dr \]
\note{« premières » est écrit au-dessus de « trois faces ». « dernieres », souligné, est écrit en tête de ligne, devant « faces », au bord gauche ; « aux premieres » est barré en fin de proposition. « chaleurs » : ainsi, au pluriel.}

le produit $c.D.r^{2}\,\text{sin}\,\theta\,d\theta\,d\varphi\,dr$ exprime combien il faut de
chaleur pour élever de 0 a 1 la molécule dont le
volume est $r^{2}\,\text{sin}\,\theta\,d\theta\,d\varphi\,dr$. En divisant donc par
ce produit la nouvelle quantité de chaleur que la molécule
vient d'obtenir on aura son accroissement de temperature
On obtiendra ainsi l'équation
\note{Le $D$ du produit $c.D$ est une capitale en forme de dôme sur une barre, comme celle que le lot 28 lit D dans le brouillon sur la chaleur de la vue 549 ; ce pourrait être un $\Omega$.}

$\frac{dv}{dt} =$ on voit qu'il suffit
\struck{si effectuer les op} d'effectuer les differentiations indiquées
pour que cette équation sera identiquement la même que celle qui \struck{celle qui a été} déd
a été déduite de l'équation générale
\note{Après $\frac{dv}{dt}$, un tiret et un blanc : le second membre de l'équation n'est pas écrit. « que » est écrit au-dessus de « cette » ; « la même que celle qui » au-dessus de « celle qui a été », barré ; « déd » qui suit n'est pas barré, et la ligne suivante reprend « a été déduite ». « sera » n'est pas corrigé en « soit ». Dans la marge de gauche, écrits de bas en haut à hauteur de ces lignes, des essais de plume : « S », « Sphere » et un mot \ill{}.}

\page{650}

\note{En haut à droite, à l'encre, « 319 ». Feuillet monté sur une feuille de garde, écrit d'un seul côté ; le dos (vue 651) est blanc. Un titre en tête, puis le calcul ; au milieu de la page, le cachet de la Bibliothèque impériale. Le bas du feuillet est blanc.}

Somme des raisons inverses des deux rayons de principale
courbure
\note{« principale » : sans marque visible de pluriel ; « courbure » finit sur un trait qui peut être un « s ». Un trait est tiré sous le titre, à gauche.}

\[ \frac{1}{R} + \frac{1}{R'} = \frac{1 + \left(\frac{dz}{dy}\right)^{2}}{\left\{1 + \left(\frac{dz}{dy}\right)^{2} + \left(\frac{dz}{dx}\right)^{2}\right\}^{\frac{3}{2}}}.\frac{d^{2}z}{dx^{2}} - \frac{2.\frac{dz}{dy}.\frac{dz}{dx}}{\left\{1 + \left(\frac{dz}{dy}\right)^{2} + \left(\frac{dz}{dx}\right)^{2}\right\}^{\frac{3}{2}}}\,\frac{d^{2}z}{dx\,dy} \]
\[ + \frac{1 + \left(\frac{dz}{dx}\right)^{2}}{\left\{1 + \left(\frac{dz}{dy}\right)^{2} + \left(\frac{dz}{dx}\right)^{2}\right\}^{\frac{3}{2}}}\,\frac{d^{2}z}{dy^{2}} \]

Si on fait $x = r\,\text{sin}\,\theta\,\text{cos.}\varphi \qquad y = r\,\text{sin}\,\theta\,\text{sin}\,\varphi$ et qu'on \struck{fasse}
prenne $r$ constant pour avoir seulement le cas relatif a
la surface dela sphère on aura
\[ \frac{dz}{dx} = \frac{\text{cos.}\theta\,\text{cos.}\varphi}{r}\,\frac{dz}{d\theta} - \frac{\text{sin}\,\varphi}{r\,\text{sin}\,\theta}.\frac{dz}{d\varphi} \]
\[ \frac{dz}{dy} = \frac{\text{cos}\,\theta\,\text{sin.}\varphi}{r}.\frac{dz}{d\theta} + \frac{\text{cos}\,\varphi}{r\,\text{sin}\,\theta}.\frac{dz}{d\varphi} \]
\[ K^{3} = \left\{1 + \left(\frac{dz}{dx}\right)^{2} + \left(\frac{dz}{dy}\right)^{2}\right\}^{\frac{3}{2}} = \left(1 + \uncertain{2}\,\frac{\text{cos}^{2}\theta}{r^{2}}\left(\frac{dz}{d\theta}\right)^{2} + \frac{1}{r^{2}\,\text{sin}^{2}\theta}.\left(\frac{dz}{d\varphi}\right)^{2}\right)^{\frac{3}{2}} \]
\note{La capitale lue $K$ est repassée. Devant « cos$^{2}\theta$ », un signe en forme de « b », lu 2 avec doute ; le calcul demande 1, et rien n'indique qu'il soit barré.}

\[ \struck{\frac{d^{2}z}{dx\,dy} = (\text{sin}\,\theta\,\text{cos}\,\varphi + \text{sin}\,\varphi} \]
\note{Un premier départ, barré d'un trait, le premier membre biffé d'un gros trait oblique.}

\[ \frac{d^{2}z}{dx\,dy} = \left\{\frac{\text{sin.}\theta\,\text{cos.}\varphi}{r}\,\frac{d\theta}{dy} + \frac{\text{cos}\,\theta\,\text{sin}\,\varphi}{r}.\frac{d\varphi}{dy}\right\}\frac{dz}{d\theta} + \frac{\text{cos}\,\theta\,\text{cos}\,\varphi}{r}\left\{\frac{d^{2}z}{d\theta^{2}}\,\frac{d\theta}{dy} + \frac{d^{2}z}{d\theta\,d\varphi}\,\frac{d\varphi}{dy}\right. \]
\[ + \left\{\frac{\text{cos}\,\varphi}{r\,\text{sin}\,\theta}\,\frac{d\varphi}{dy} - \frac{\text{sin}\,\varphi\,\text{cos}\,\theta}{r\,\text{sin}^{2}\theta}.\frac{d\theta}{dy}\right\}\frac{dz}{d\varphi} - \frac{\text{sin}\,\varphi}{r\,\text{sin.}\theta}\left\{\frac{d^{2}z}{d\varphi\,d\theta}\,\frac{d\theta}{dy} + \frac{d^{2}z}{d\varphi^{2}}\,\frac{d\varphi}{dy}\right. \]
\note{Aucun signe n'est écrit devant le premier terme de la première accolade. Les accolades de droite ne se ferment pas : le dernier facteur touche le bord du feuillet. Le $\varphi$ de « cos$\,\varphi$ », dans la première accolade, est récrit.}

\[ = \left\{\frac{\text{sin}\,\theta\,\text{cos}\,\theta\,\text{sin}\,\varphi\,\text{cos}\,\varphi}{r^{2}} + \frac{\text{sin}\,\theta\,\text{cos}\,\theta\,\text{sin}\,\varphi\,\text{cos}\,\varphi}{r\,\text{sin}^{2}\theta}\right\}\frac{dz}{d\theta} + \frac{\text{cos}^{2}\theta\,\text{cos}^{2}\varphi}{r^{2}}\,\frac{d^{2}z}{d\theta^{2}} \]
\[ + \left\{\frac{\text{cos}^{2}\varphi}{r^{2}\,\text{sin}^{2}\theta} - \frac{\text{cos}^{2}\theta\,\text{sin}^{2}\varphi}{r^{2}\,\text{sin}^{2}\theta}\right\}\frac{dz}{d\varphi} + \frac{\text{cos}\,\theta\,(\text{cos}^{2}\varphi - \text{sin}^{2}\varphi)}{r^{2}\,\text{sin}\,\theta}.\frac{d^{2}z}{d\varphi\,d\theta} - \frac{\text{cos}\,\varphi\,\text{sin}\,\varphi}{r^{2}\,\text{sin}^{2}\theta}\,\frac{d^{2}z}{d\varphi^{2}} \]
\note{Le dénominateur du deuxième terme de la première accolade est lu « $r\,\text{sin}^{2}\theta$ », sans exposant à $r$. L'exposant de « cos$^{2}\theta$ » au dernier terme de la première ligne et celui de « sin$^{2}\varphi$ » dans la seconde sont récrits.}

\page{652}

\note{En haut à droite, à l'encre, « 320 » ; au-dessus, d'un trait plus épais et plus haut, un « 1 » suivi d'une petite croix, dont le sens n'est pas établi (une première page ? un renvoi ?). Feuillet monté sur une feuille de garde, écrit d'un seul côté, d'une écriture régulière ; le dos (vue 653) est blanc. Cachet de la Bibliothèque impériale en marge droite. Le bas du feuillet est blanc, avec des taches au bord inférieur.}

Équation du mouvement varié dela chaleur dans une sphère
solide échauffée d'une maniere entièrement arbitraire
\note{Titre, suivi d'un court trait centré.}

Pour appliquer l'équation générale $\frac{dv}{dt} = \frac{K}{c.D}\left(\frac{d^{2}v}{dx^{2}} + \frac{d^{2}v}{dy^{2}} + \frac{d^{2}v}{dz^{2}}\right)$
au cas d'une figure sphérique il faut exprimer les coordonnées
en fonction du rayon et des \struck{angles qu'une position donnée
de cette ligne fait} quantités angulaires
\note{Le $D$ de « $c.D$ » est la même capitale en dôme qu'à la vue 646.}

Prenons \struck{l'origine au} centre de la sphère \struck{pour} origine
en adoptant les dénominations qui appartiennent a la théorie
dela terre nous prendrons le plan de l'équateur pour
celui des $x$ et $y$ parconséquent la valeur complette
de $z$ sera égal au rayon dela sphère et elle appartiendra
au pôle de ce solide.
\note{La correction de la première ligne est laissée incomplète : barrés « l'origine au » et « pour », elle se lit « Prenons centre de la sphère origine ». « égal » : sans accord.}

Soit $\theta$ l'arc compris entre le pôle et un point
quelconque dela sphère on aura pour ce point
$z = r\,\text{cos}\,\theta \qquad r\,\text{sin}\,\theta$ sera le rayon du cercle formé par
l'intersection du plan mené par le même point parallelement
au plan de l'équateur. si $\varphi$ est l'arc compris entre
le point dont il s'agit et celui auquel appartient
la plus grande valeur de l'ordonnée $x$ dans le parallele
auquel le point appartient on aura
\[ x = r\,\text{sin}\,\theta\,\text{cos}\,\varphi \qquad y = r\,\text{sin}\,\theta\,\text{sin.}\,\varphi. \]
\[ dx = dr\,\text{sin}\,\theta\,\text{cos}\,\varphi + r\,\text{cos}\,\theta\,\text{cos}\,\varphi\,d\theta - r\,\text{sin}\,\theta\,\text{sin}\,\varphi\,d\varphi \]
\[ dy = dr\,\text{sin.}\theta\,\text{sin.}\varphi + r\,\text{cos}\,\theta\,\text{sin}\,\varphi\,d\theta + r\,\text{sin}\,\theta\,\text{cos}\,\varphi\,d\varphi \]
\[ dz = dr\,\text{cos}\,\theta - r\,\text{sin}\,\theta\,d\theta \]
de ces valeurs de $dx$ $dy$ et $dz$ on tire
\note{« parallelement » finit serré contre le bord droit. Dans $y$, « sin.$\,\varphi$ » est écrit sur un premier mot noirci. Les deux signes $-$ de $dx$ et de $dz$ sont des traits empâtés en forme de tache. La page s'arrête sur « on tire » ; la suite n'est pas au dos (vue 653, blanc) : elle peut être le calcul du feuillet 321 (vue 654), qui commence par les mêmes combinaisons de $dx$, $dy$, $dz$.}

\page{654}

\note{En haut à droite, à l'encre, « 321 » ; au-dessus, d'un trait plus épais, un chiffre barré d'un trait oblique, lu 2 avec doute (ce pourrait être un 8), qui fait suite au « 1 » marqué au-dessus de « 320 » à la vue 652. Feuillet monté sur une feuille de garde, écrit d'un seul côté, d'une écriture régulière ; le dos (vue 655) est blanc. Calcul sans texte : il fait suite à « de ces valeurs de $dx$ $dy$ et $dz$ on tire » de la vue 652.}

\[ dx.\,\text{sin}\,\theta\,\text{cos}\,\varphi + dy\,\text{sin.}\theta\,\text{sin.}\varphi + dz\,\text{cos.}\theta \]
\[ = dr\,(\text{sin}^{2}\theta\,\text{cos}^{2}\varphi + \text{sin}^{2}\theta\,\text{sin}^{2}\varphi + \text{cos}^{2}\theta) + r\,d\theta\,(\text{cos}\,\theta\,\text{sin}\,\theta\,(\text{cos}^{2}\varphi + \text{sin}^{2}\varphi) - 1) \]
\[ + r\,d\varphi\,(\text{sin}^{2}\theta\,(\text{sin}\,\varphi\,\text{cos}\,\varphi - \text{sin}\,\varphi\,\text{cos}\,\varphi)) \]
\[ = dr \]
\note{Dans le coefficient de $r\,d\theta$, « $-1$ » est écrit tel quel, là où l'on attend $-\text{cos}\,\theta\,\text{sin}\,\theta$ ; la parenthèse avant « $-1$ » est repassée.}

\[ dx.\,\text{cos.}\theta\,\text{cos.}\varphi + dy\,\text{cos}\,\theta\,\text{sin}\,\varphi - dz.\,\text{sin}\,\theta \]
\[ = dr\,(\text{cos.}\theta\,\text{sin}\,\theta\,(\text{cos}^{2}\varphi + \text{sin}^{2}\varphi - 1)) + r\,d\theta\,(\text{cos}^{2}\theta\,\text{cos}^{2}\varphi + \text{cos}^{2}\theta\,\text{sin}^{2}\varphi + \text{sin}^{2}\theta) \]
\[ + r\,d\varphi\,(\text{sin}\,\theta\,\text{cos}\,\theta\,(\text{sin}\,\varphi\,\text{cos}\,\varphi - \text{sin}\,\varphi\,\text{cos}\,\varphi)) \]
\[ = r\,d\theta \]

\[ \struck{dx\,\text{cos}\ \ill{}}\ dy\,\text{cos}\,\varphi - dx\,\text{sin.}\varphi \]
\[ = dr\,\text{sin}\,\theta\,(\text{sin.}\varphi\,\text{cos}\,\varphi - \text{sin.}\varphi\,\text{cos}\,\varphi) + r\,\text{cos}\,\theta\,d\theta\,(\text{sin}\,\varphi\,\text{cos}\,\varphi - \text{sin}\,\varphi\,\text{cos}\,\varphi) \]
\[ + r\,\text{sin}\,\theta\,d\varphi\,(\text{cos}^{2}\varphi + \text{sin}^{2}\varphi) \]
\[ = r\,\text{sin}\,\theta\,d\varphi \]
\note{La ligne commence par un premier essai biffé de hachures serrées, où se lisent « $dx$ cos » ; les « cos » et « sin » qui suivent $dy$ et $dx$ sont récrits sur des mots noircis.}

\[ \frac{dr}{dx} = \text{sin.}\theta\,\text{cos.}\varphi \qquad \frac{d\theta}{dx} = \frac{\text{cos.}\theta\,\text{cos}\,\varphi}{r}, \qquad \frac{d\varphi}{dx} = -\frac{\text{sin}\,\varphi}{r\,\text{sin.}\theta} \]
\[ \frac{dr}{dy} = \text{sin.}\theta\,\text{sin.}\varphi \qquad \frac{d\theta}{dy} = \frac{\text{cos.}\theta\,\text{sin.}\varphi}{r}, \qquad \frac{d\varphi}{dy} = \frac{\text{cos}\,\varphi}{r\,\text{sin.}\theta} \]
\[ \frac{dr}{dz} = \text{cos}\,\theta \qquad \frac{d\theta}{dz} = -\frac{\text{sin.}\theta}{r}, \]
\note{Le « sin » de $\frac{dr}{dy}$ est récrit sur un autre mot.}

\[ \frac{dv}{dx} = \frac{dv}{dr}.\frac{dr}{dx} + \frac{dv}{d\theta}.\frac{d\theta}{dx} + \frac{dv}{d\varphi}.\frac{d\varphi}{dx} \]
\[ = \frac{dv}{dr}.\,\text{sin.}\theta\,\text{cos.}\varphi + \frac{dv}{d\theta}.\frac{\text{cos.}\theta\,\text{cos.}\varphi}{r} - \frac{dv}{d\varphi}\,\frac{\text{sin.}\varphi}{r\,\text{sin}\,\theta} \]
\[ \frac{dv}{dy} = \frac{dv}{dr}.\frac{dr}{dy} + \frac{dv}{d\theta}.\frac{d\theta}{dy} + \frac{dv}{d\varphi}.\frac{d\varphi}{dy} \]
\[ = \frac{dv}{dr}.\,\text{sin.}\theta\,\text{sin}\,\varphi + \frac{dv}{d\theta}.\frac{\text{cos.}\theta\,\text{sin.}\varphi}{r} + \frac{dv}{d\varphi}\,\frac{\text{cos}\,\varphi}{r\,\text{sin.}\theta} \]
\[ \frac{dv}{dz} = \frac{dv}{dr}.\frac{dr}{dz} + \frac{dv}{d\theta}\,\frac{d\theta}{dz} \]
\[ = \frac{dv}{dr}\,\text{cos}\,\theta - \frac{dv}{d\theta}\,\frac{\text{sin}\,\theta}{r} \]
\note{Le $v$ du numérateur $\frac{dv}{d\varphi}$, dans l'expression développée de $\frac{dv}{dx}$, est récrit sur une tache. Dans les lignes développées, le $v$ de $\frac{dv}{d\theta}$ porte une queue descendante qui le fait ressembler à un « y » (de même à la vue 643) ; il est lu $v$ d'après les lignes générales. La page s'arrête là ; le bas du feuillet est blanc.}

\page{656}

\note{En haut à droite, à l'encre, « 322 », souligné d'un trait ; à gauche du nombre, un peu plus haut, le mot « Note », d'une plume fine : la page se donne comme une note. Feuillet monté sur une feuille de garde, écrit d'un seul côté, d'une écriture régulière corrigée entre les lignes ; le dos (vue 657) est blanc. Le bas du feuillet est blanc.}

Si dans l'équation (A) p. 3 dela sphère échauffée d'une
maniere arbitraire on fait $r = \infty \quad \text{cos}\,\theta = \ill{} \quad \text{sin}\,\theta = \theta \quad r\,\text{sin}\,\theta$ sera
une quantité finie qu'on peut représenter par \emph{s}.
\note{Au-dessus de « $r\,\text{sin}\,\theta$ sera », dans l'interligne, en deux lignes serrées : « $\theta$ étant infiniment petit ». La valeur de $\text{cos}\,\theta$ est couverte d'une tache. Le « s » de « par s » est souligné. L'équation marquée (A) est celle qui clôt le feuillet 324 (vue 660), relié après cette note ; « p. 3 » ne répond à aucun chiffre lu sur ces feuillets (« 1 » au-dessus de 320, un « 2 » douteux au-dessus de 321).}

Cela posé on aura $\frac{dv}{dt} = \frac{K}{c.D}\left\{\frac{d^{2}v}{dr^{2}} + \frac{d^{2}v}{ds^{2}} + \frac{1}{s}.\frac{dv}{ds} + \frac{d^{2}v}{s^{2}d\varphi^{2}}\right\}$

Car le terme $\frac{2}{r}\,\frac{dv}{dr}$ devient nul puis qu'il a $\frac{2}{\infty}$ pour facteur
cette équation est celle du cylindre car $\frac{d^{2}v}{dr^{2}}$ est la même
chose que \uncertain{$x$}. cette maniere dela déduire est plus simple
que celle imployée par M\textsuperscript{r} Poisson car je ne vois pas
la necessité de mettre $r - a$ a la place de $r$. si ce n'est
peut-être pour mieux faire sentir que les ordonnées \struck{s}ont
le point où $x = 0$ pour origine mais cette substitution
peut avoir lieu après coup elle expliquera que la variation
de $r$ supposé infini ne peut s'appliquer qu'a la \uncertain{ba}
quantité finie $x$ qui lui est \struck{ajoute} ajouté lorsqu'on
fait $\uncertain{c} = r + x$
\note{« que $x$ » : une seule lettre, lue $x$ avec doute, suivie d'un point. « imployée » : ainsi. « ont » est récrit sur une première lettre noircie. « substitution » est récrit sur un premier mot ; « le point » est écrit en un mot. En fin de ligne, après « qu'a la », deux lettres, lues « ba », sans suite : un mot commencé et laissé ; la ligne suivante reprend « quantité finie $x$ ». La lettre lue $c$ dans « $c = r + x$ » est nette mais surprend ; ce pourrait être un $r$ mal formé.}

Si dans la même équation (A) on \struck{suppose qu'} substitue a
$\frac{dv}{d\theta} \quad \frac{dv}{d(\text{sin}\,\theta)}$ elle deviendra
\[ \frac{dv}{dt} = \frac{K}{C.D}\left\{\frac{2}{r}\,\frac{dv}{dr} + \frac{d^{2}v}{dr^{2}} + \frac{d^{2}v}{r^{2}\,\text{cos}^{2}\theta\,d\theta^{2}} + \frac{\text{cos}\,\theta}{r\,\text{sin}\,\theta}.\frac{dv}{r\,\text{cos}\,\theta\,d\theta} + \frac{1}{r^{2}\,\text{sin}^{2}\theta}.\frac{d^{2}v}{d\varphi^{2}}\right\} \]
\note{Tout ce paragraphe et son équation sont annulés par deux longs traits croisés en X, qui vont de « Si dans la même » à la fin de l'équation. Le mot « même » porte un petit trait au-dessus. Le $r$ du dénominateur de $\frac{dv}{dr}$ est récrit ; dans le dernier terme, le dénominateur « $r^{2}\,\text{sin}^{2}\theta$ » est récrit sur une tache.}

On s'étonneroit sans doute de voir l'équation du cylindre comprise
dans celle dela sphère si on ne voyoit qu'en cherchant cette
derniere d'une maniere spéciale on est conduit a considérer le solide
conique qui s'appuie sur la calotte spérique dont la base seroit
en même tems celle du cône. On concoit que lorsque le rayon
est infini la calotte spérique se confond avec sa base et
que parconséquent le solide conique qui a fourni l'équation
de la sphère devient un cylindre.
\note{« spérique » : ainsi, deux fois, sans h. « concoit » sans cédille. Le texte s'arrête là, au milieu du feuillet.}

\page{660}

\note{En haut à droite, à l'encre, « 324 », un trait oblique tiré sous le nombre ; aucune autre marque en tête. Feuillet monté sur une feuille de garde, écrit d'un seul côté, d'une écriture régulière ; le dos (vue 661, hors de ce lot) est blanc au vu de la vignette. Cachet de la Bibliothèque impériale dans la marge droite. Le calcul continue celui du feuillet 321 (vue 654), qui s'arrêtait sur $\frac{dv}{dx}$, $\frac{dv}{dy}$, $\frac{dv}{dz}$, et finit sur l'équation marquée (A).}

\[ \frac{d^{2}v}{dx^{2}} = \text{sin.}\theta\,\text{cos.}\varphi\left\{\frac{d^{2}v}{dr^{2}}.\frac{dr}{dx} + \frac{d^{2}v}{dr\,d\theta}.\frac{d\theta}{dx} + \frac{d^{2}v}{dr\,d\varphi}.\frac{d\varphi}{dx}\right\} \]
\[ + \frac{dv}{dr}\left\{\text{cos}\,\theta\,\text{cos}\,\varphi\,\frac{d\theta}{dx} - \text{sin}\,\theta\,\text{sin.}\varphi\,\frac{d\varphi}{dx}\right\} \]
\[ + \frac{\text{cos.}\theta\,\text{cos.}\varphi}{r}\left\{\frac{d^{2}v}{d\theta\,dr}.\frac{dr}{dx} + \frac{d^{2}v}{d\theta\,d\varphi}.\frac{d\varphi}{dx} + \frac{d^{2}v}{d\theta^{2}}.\frac{d\theta}{dx}\right\} \]
\[ - \frac{dv}{d\theta}\left\{\frac{\text{sin}\,\theta\,\text{cos}\,\varphi}{r}\,\frac{d\theta}{dx} + \frac{\text{sin}\,\varphi\,\text{cos}\,\theta}{r}\,\frac{d\varphi}{dx} - \frac{\text{cos}\,\theta\,\text{cos}\,\varphi}{r^{2}}\,\frac{dr}{dx}\right\} \]
\[ - \frac{\text{sin.}\varphi}{r\,\text{sin.}\theta}\left\{\frac{d^{2}v}{d\varphi\,dr}.\frac{dr}{dx} + \frac{d^{2}v}{d\varphi\,d\theta}\,\frac{d\theta}{dx} + \frac{d^{2}v}{d\varphi^{2}}\,\frac{d\varphi}{dx}\right\} \]
\[ - \frac{dv}{d\varphi}\left\{\frac{\text{cos}\,\varphi}{r\,\text{sin}\,\theta}\,\frac{d\varphi}{dx} - \frac{\text{sin.}\varphi\,\text{cos}\,\theta}{r\,\text{sin}^{2}\theta}\,\frac{d\theta}{dx} - \frac{\text{sin.}\varphi}{r^{2}\,\text{sin.}\theta}\,\frac{dr}{dx}\right\} \]
\note{Dans l'accolade de $\frac{dv}{d\theta}$, le $\varphi$ de « cos$\,\varphi$ » est récrit sur une autre lettre.}

\[ = \frac{d^{2}v}{dr^{2}}.\,\text{sin}^{2}\theta\,\text{cos}^{2}\varphi + \frac{d^{2}v}{dr\,d\theta}.\frac{2\,\text{sin}\,\theta\,\text{cos}\,\theta\,\text{cos}^{2}\varphi}{r} - \frac{d^{2}v}{dr\,d\varphi}.\frac{2\,\text{sin}\,\varphi\,\text{cos}\,\varphi}{r} \]
\[ + \frac{d^{2}v}{d\theta^{2}}\,\frac{\text{cos}^{2}\theta\,\text{cos}^{2}\varphi}{r^{2}} - \frac{d^{2}v}{d\theta\,d\varphi}\,\frac{2\,\text{cos}\,\theta\,\text{sin}\,\varphi\,\text{cos}\,\varphi}{r^{2}\,\text{sin.}\theta} + \frac{d^{2}v}{d\varphi^{2}}\,\frac{\text{sin}^{2}\varphi}{r^{2}\,\text{sin}^{2}\theta} \]
\[ + \frac{dv}{dr}\left\{\frac{\text{cos}^{2}\theta\,\text{cos}^{2}\varphi}{r} + \frac{\text{sin}^{2}\varphi}{r}\right\} - \frac{dv}{d\theta}\left\{\frac{2\,\text{sin.}\theta\,\text{cos}\,\theta\,\text{cos}^{2}\varphi}{r^{2}} - \frac{\text{sin}^{2}\varphi\,\text{cos}\,\theta}{r^{2}\,\text{sin}\,\theta}\right\} \]
\[ + \frac{dv}{d\varphi}\left\{\frac{\text{sin}\,\varphi\,\text{cos}\,\varphi}{r^{2}\,\text{sin}^{2}\theta} + \frac{\text{cos}^{2}\theta\,\text{sin}\,\varphi\,\text{cos}\,\varphi}{r^{2}\,\text{sin}^{2}\theta} + \frac{\text{sin}\,\varphi\,\text{cos}\,\varphi}{r^{2}}\right\} \]
\note{Le $v$ de $\frac{d^{2}v}{dr\,d\theta}$ et celui de $\frac{d^{2}v}{d\theta\,d\varphi}$ sont récrits sur des taches ; le signe $-$ devant ce dernier est un trait empâté. Après l'accolade de $\frac{dv}{dr}$, un trait vertical barré ; l'accolade de $\frac{dv}{d\theta}$ s'ouvre sur un signe récrit. Au dernier dénominateur de cette accolade, une petite tache à la place d'un éventuel exposant de « sin$\,\theta$ » ; lu sans exposant.}

\[ \frac{d^{2}v}{dy^{2}} = \frac{d^{2}v}{dr^{2}}\,\text{sin}^{2}\theta\,\text{sin}^{2}\varphi + \frac{d^{2}v}{dr\,d\theta}.\frac{2\,\text{sin}\,\theta\,\text{cos}\,\theta\,\text{sin}^{2}\varphi}{r} + \frac{d^{2}v}{dr\,d\varphi}.\frac{2\,\text{sin}\,\varphi\,\text{cos}\,\varphi}{r} \]
\[ + \frac{d^{2}v}{d\theta^{2}}\,\frac{\text{cos.}^{2}\theta\,\text{sin.}^{2}\varphi}{r^{2}} + \frac{d^{2}v}{d\theta\,d\varphi}.\frac{2\,\text{cos}\,\theta\,\text{sin.}\varphi\,\text{cos}\,\varphi}{r^{2}\,\text{sin.}\theta} + \frac{d^{2}v}{d\varphi^{2}}\,\frac{\text{cos}^{2}\varphi}{r^{2}\,\text{sin.}^{2}\theta} \]
\[ + \frac{dv}{dr}\left\{\frac{\text{cos.}^{2}\theta\,\text{sin}^{2}\varphi}{r} + \frac{\text{cos}^{2}\varphi}{r}\right\} - \frac{dv}{d\theta}\left\{\frac{2\,\text{sin}\,\theta\,\text{cos}\,\theta\,\text{sin}^{2}\varphi}{r^{2}} - \frac{\text{cos}^{2}\varphi\ \struck{\text{cos}\,\theta}}{r^{2}\,\text{sin}^{\ill{}}\theta}\right) \]
\[ - \frac{dv}{d\varphi}\left\{\frac{\text{sin}\,\varphi\,\text{cos}\,\varphi}{r^{2}\,\text{sin}^{2}\theta} + \frac{\text{cos}^{2}\theta\,\text{sin}\,\varphi\,\text{cos}\,\varphi}{r^{2}\,\text{sin}^{2}\theta} + \frac{\text{sin}\,\varphi\,\text{cos}\,\varphi}{r^{2}}\right\} \]
\note{Le signe devant $\frac{dv}{d\theta}$ est une tache ; lu $-$. L'exposant de « sin$\,\theta$ » au dernier dénominateur de cette accolade est récrit et ne se lit pas ; l'accolade se ferme par une parenthèse.}

\[ \frac{d^{2}v}{dx^{2}} + \frac{d^{2}v}{dy^{2}} = \text{sin}^{2}\theta.\,\frac{d^{2}v}{dr^{2}} + \frac{2\,\text{sin.}\theta\,\text{cos.}\theta}{r}\,\frac{d^{2}v}{dr\,d\theta} + \frac{\text{cos}^{2}\theta}{r^{2}}.\frac{d^{2}v}{d\theta^{2}} + \frac{1}{r^{2}\,\text{sin}^{2}\theta}\,\frac{d^{2}v}{d\varphi^{2}} \]
\[ + \frac{dv}{dr}\left\{\frac{\text{cos}^{2}\theta}{r} + \frac{1}{r}\right\} - \frac{dv}{d\theta}\left\{\frac{2\,\text{sin.}\theta\,\text{cos.}\theta}{r^{2}} - \frac{\text{cos.}\theta}{r^{2}\,\text{sin.}\theta}\right\} \]
\[ \frac{d^{2}v}{dz^{2}} = \text{cos}^{2}\theta\,\frac{d^{2}v}{dr^{2}} - \frac{2\,\text{sin}\,\theta\,\text{cos.}\theta}{r}\,\frac{d^{2}v}{dr\,d\theta} + \frac{\text{sin}^{2}\theta}{r^{2}}\,\frac{d^{2}v}{d\theta^{2}} \]
\[ + \frac{dv}{dr}.\frac{\text{sin}^{2}\theta}{r} + \frac{dv}{d\theta}\,\frac{\uncertain{2}\,\text{sin.}\theta\,\text{cos.}\theta}{r^{2}} \]
\note{Le coefficient du dernier terme est écrit sur une tache, lu 2 avec doute.}

\[ \frac{d^{2}v}{dx^{2}} + \frac{d^{2}v}{dy^{2}} + \frac{d^{2}v}{dz^{2}} = \frac{2}{r}.\frac{dv}{dr} + \frac{d^{2}v}{dr^{2}} + \frac{d^{2}v}{r^{2}d\theta^{2}} + \frac{\text{cos}\,\theta}{r^{2}\,\text{sin}\,\theta}.\frac{dv}{d\theta} + \frac{\uncertain{1}}{r^{2}\,\text{sin}^{2}\theta}\,\frac{d^{2}v}{d\varphi^{2}} \]
\note{Le numérateur du dernier terme est une tache, lue 1 avec doute d'après l'équation qui suit.}

L'équation générale devient donc
\[ \frac{dv}{dt} = \frac{K}{c.D}\left\{\frac{2}{r}\,\frac{dv}{dr} + \frac{d^{2}v}{dr^{2}} + \frac{d^{2}v}{r^{2}d\theta^{2}} + \frac{\text{cos.}\theta}{r^{2}\,\text{sin.}\theta}.\frac{dv}{d\theta} + \frac{1}{r^{2}\,\text{sin}^{2}\theta}\,\frac{d^{2}v}{d\varphi^{2}}\right\} \]
\[ \ldots\ (\uncertain{\text{A}}) \]
\note{La lettre de renvoi, entre parenthèses après quatre points, est récrite et empâtée ; la lecture (A) s'accorde avec « l'équation (A) … dela sphère échauffée d'une maniere arbitraire » de la vue 656. Le bas du feuillet est blanc.}

\end{document}
